\section{The Post-Merge Validation Loop}
\label{sec:validation}

% ~1.5 pages. This section carries the thesis empirically. Do not cut it.
% STATUS: [spec] -- lands with merge policies and the tracker connectors.

\subsection{Intent validation is above the gaming layer}
% Gates verify correctness against specification; they are attackable by the
% thing being gated. A human confirming "this does what the ticket asked" is not
% reading the diff -- they are testing behaviour against the requirement.
% Suppression comments and weakened assertions do not survive that check.
% Crucially it is O(tickets), not O(agent-turns): a task that took 200 turns and
% 3 candidates gets exactly one validation.

\subsection{Failure is a normal outcome}
% Post-merge validation failure must be as cheap and structured as a gate
% failure: revert, new task carrying the finding, and -- the important part --
% the finding becomes a gate so it cannot recur. Concretely: a new gate class,
% made mandatory through the cluster policy's required gates, injected into
% every plan admitted from then on (Section 7.5).

\subsection{Compounding trust}
% This loop is the mechanism by which the system becomes more trustworthy over
% time. It is also what justifies raising a repo's position on the autonomy dial.

\subsection{Sampled audit}
% Audit k\% of auto-merged PRs; track the rate at which audits catch what gates
% missed. Constant cost, not linear in task count. This is the empirical trust
% measure and the basis for the mergePolicy decision.
